%% Elemento obrigatório (Figura 15).
%% Consiste em uma tradução do resumo em português para uma
%% língua estrangeira (em inglês, ABSTRACT; em espanhol, RESUMEN;
%% em francês, RÉSUMÉ), em um único parágrafo, seguido das palavras-
%% -chave representativas do conteúdo do trabalho, na língua estrangeira
%% escolhida.
%% O resumo em outra língua também é precedido pela referência
%% do trabalho, substituindo-se o título em português pelo título na língua
%% estrangeira adotada.
%% No caso de teses, é possível incluir dois resumos em língua es-
%% trangeira.
%% A apresentação gráfica e a ordem dos elementos seguem a mes-
%% ma orientação do resumo em português.

\begin{resumo}[Abstract]
\begin{otherlanguage*}{english}
\begin{SingleSpace}

\noindent
\begin{flushleft}
\entradaAutor{}. \textit{\englishTitle{}}. \the\year. \pageref{LastPage} f. Trabalho de Conclusão de Curso (Graduação em Engenharia de Computação) - Instituto Politécnico, Universidade do Estado do Rio de Janeiro, Nova Friburgo, \the\year.
\end{flushleft}
\vspace{\onelineskip}

\setlength{\parindent}{1.3cm}

This work presents the development of Onmarmoraria, a multi-tenant ERP system for marble businesses, aimed at centralizing clients, products, and quotes in a web dashboard. The methodology included requirements gathering, data modeling, implementation of a Go REST API with JWT authentication, a Next.js frontend, and automated plus manual testing. The infrastructure uses Google Cloud Platform (Compute Engine and Cloud Storage), Docker, GHCR, GitHub Actions, and Vercel, prioritizing low operational cost and continuous delivery. Application media files are stored in GCS via signed URLs, while Docker images of the API are published to GHCR.

\vspace{\onelineskip}
\noindent Keywords: ERP; Marble Industry; Quotes; Cloud Computing; Google Cloud Platform.

\end{SingleSpace}
\end{otherlanguage*}
\end{resumo}